We have presented the Filtered Stepper Calculus, a formal framework
for scoped, pattern-based control over which reduction steps are
visible during algebraic stepping. The calculus is backed by
mechanized proofs of preservation, progress, and simulation in Agda,
and is implemented in the Hazel live programming environment, where it
has been deployed in a university programming languages course.

\paragraph{Stepper adoption.}
Our evaluation shows that students adopted the stepper organically:
without being instructed to use it, nearly half of the students who
observed the stepper in lecture chose to use it on their assignments.
This suggests that an algebraic stepper that mirrors the
substitution-based reasoning taught in class fills a genuine need
in the pedagogical workflow.

\paragraph{Filter adoption.}
By contrast, we observed limited student use of the filter system
during assignments. We attribute this to three factors. First, the
current interface requires users to write filter annotations directly
in their program text, which is unfamiliar syntax and adds friction.
Second, inserting a filter currently triggers re-elaboration of the
entire program, introducing noticeable latency. Third, students
working on assignments are focused on correctness rather than trace
exploration, so the cost of learning the filter language may not be
justified in that context.

\paragraph{Lecture versus assignment contexts.}
A natural reading is that filters answer different needs in lecture
versus assignment contexts. In lecture, the instructor prepares
filters ahead of time to direct attention, and the cost of writing
them is amortized over a classroom of viewers. On assignments, a
student writes filters for an audience of one and pays the full cost
themselves, with no immediate payoff in terms of grade or
correctness. The lecture-side adoption suggests filters are most
useful, at least initially, as a presentation tool. Broader student
adoption may depend on lowering the per-use cost (cf. the graphical
panel discussed below).

\paragraph{Limitations.}
The primary contribution of this paper is a theoretical treatment of
pattern-based control over algebraic stepping: a formal calculus,
mechanized metatheory, and a working implementation. The classroom
deployment we report is illustrative rather than a controlled study.
It involves one instructor at one institution, the assignment sample
is small ($n = 33 + 20$), survey responses are voluntary and
self-reported, and we did not measure learning outcomes. Stronger
claims about pedagogical effectiveness would require a controlled
study against an unfiltered stepper and a no-stepper baseline,
ideally across multiple instructors. Our evaluation is intended
only to demonstrate that students and an instructor will pick up the
tool and use it under normal classroom conditions.

\paragraph{Future work.}

The filter-adoption findings point to a clear design direction:
separating the filter interface from the program text. We envision a
graphical filter panel in which users can add, modify, and remove
filters interactively \emph{during stepping}, without editing source
code or triggering re-evaluation. This would lower the barrier to
entry and, perhaps more importantly, make filters useful for
\emph{interactive debugging}: a user could discover what to focus on
while exploring a trace, rather than predicting it ahead of time.
Filter annotations would then function less like up-front program
declarations and more like queries on an ongoing computation.

A second direction is a more expressive filter language. The current
pattern language matches on the syntactic structure of the redex up
to a fixed nesting depth. Two kinds of queries we found ourselves
wanting are not expressible. The first refers to the type or shape
of sub-expressions: ``skip all \texttt{let} reductions whose
right-hand side is a function value'' is a natural request, but our
patterns can only match values uniformly, not values of a particular
type. The second requires recursive structure: ``stop at any
expression composed entirely of literal values and arithmetic
operators, before it begins reducing'' would require something like
a grammar-style production or a recursive predicate over the AST.
Type-based patterns and predicate-based patterns would each
address part of this gap.

A third direction is filter inference. When a user is overwhelmed by
a long trace, the system could observe which steps they advance
through quickly versus dwell on and propose filter annotations
automatically. This would lower the activation cost for occasional
users for whom writing a filter from scratch is too much overhead.

Relatedly, instructors could ship pedagogical filter sets with
assignments---a shared library of common filter patterns that
students apply directly. Filters become reusable classroom artifacts
in their own right, and standard filters could be developed and
refined across cohorts and courses.

Finally, adapting the filter calculus to richer evaluation
strategies, such as lazy evaluation or concurrent reduction, would
broaden its applicability beyond Hazel.

\paragraph{Conclusion.}
The Filtered Stepper Calculus offers a principled, mechanized answer
to the ``too many steps'' problem in algebraic stepping: let the user
describe, lexically and by structure, which reductions to surface,
and let the calculus handle the rest. The remaining design space
lies largely in user interface---how best to expose the calculus's
expressive power to instructors and students in practice.
